% Replacement experimental section for the QSD-DETR manuscript.
% Required packages: amsmath, amssymb, array, booktabs, graphicx, multirow.
% Place the replacement figures at the first paths specified below.
% Existing artwork is displayed only as an explicitly marked temporary preview.
% The four method figures plus the introductory figure precede this section.
% Consequently, the first experimental figure is automatically Fig. 6.

% Keep the same font in every table; do not resize individual tables.
\providecommand{\QSDexpTableStyle}{%
  \footnotesize
  \setlength{\tabcolsep}{3pt}%
  \renewcommand{\arraystretch}{1.12}%
}

% #1 = final artwork, #2 = existing temporary artwork, #3 = panel height.
% The temporary notice disappears automatically when #1 is supplied.
\providecommand{\QSDexpArtwork}[3]{%
  \IfFileExists{#1}{%
    \includegraphics[width=\textwidth]{#1}%
  }{%
    \IfFileExists{#2}{%
      \includegraphics[width=\textwidth,height=#3,keepaspectratio]{#2}\par
    }{%
      \fbox{\parbox[c][#3][c]{0.94\textwidth}{\centering
        Experimental figure reserved for replacement artwork.}}\par
    }%
    \smallskip
    {\footnotesize\fbox{\strut Temporary artwork: replace with the figure described in the caption.}}%
  }%
}

\section{Experimental Results}
\label{sec:experiments}

This section evaluates whether QSD-DETR retains known-component detection
accuracy while providing useful image-level evidence for unknown defects.
The evaluation separates component localization, unknown-defect ranking,
training-objective ablation, and empirical selective rejection. This
separation is necessary because accurate component detection does not imply
that detector confidence can identify a previously unseen physical defect.
The proposed configurations are denoted \emph{Ours} in tables and figures;
RT-DETR denotes the detector baseline, instantiated with RT-DETRv2 and a
ResNet-50 backbone.

\subsection{Experimental Protocol and Metrics}
\label{subsec:experimental_protocol}

\subsubsection{Known-Component Detection}
The closed-set comparison uses the InsPLAD validation split containing
2626 images and 6324 annotated instances from 17 component categories.
RT-DETR and QSD-DETR use a checkpoint-compatible annotation file whose
category identifiers run from 1 to 18. Its additional \emph{sphere} entry
contains no validation instances and has no valid category precision in the
COCO evaluation cache; it is excluded from the reported category average.
The 17-category annotation file used for the YOLO reports contains the same
validation image and instance counts, with identifiers remapped to 0--16.
The compared detectors are YOLOv8m, YOLO11m, YOLOv12m, Faster R-CNN with a
ResNet-50-FPN backbone, RT-DETR, and QSD-DETR.

Detection results are expressed as percentages. We report mAP$_{50}$ and
mAP$_{75}$ at intersection-over-union (IoU) thresholds of 0.50 and 0.75,
respectively, and mAP$_{50:95}$ averaged over ten thresholds from 0.50 to
0.95 in steps of 0.05. The additional mAP$_{50:75}$ averages six thresholds
from 0.50 to 0.75 in the same steps; it is distinct from mAP$_{75}$.
Unavailable entries in the archived reports are marked ``--'' rather than
estimated from another metric. The RT-DETR results are recovered from the
EMA evaluation cache, whereas the QSD-DETR results are recovered from the
non-EMA cache. Their difference therefore describes the evaluated
checkpoints and is not an isolated causal estimate of an auxiliary module.

\subsubsection{Image-Level Unknown-Defect Assessment}
The condition-controlled comparison evaluates rain, snow, and inpainting
separately. Each condition contains 422 test images: 259 normal images and
163 images with unknown defects. The saved method comparisons use the same
image identities within a condition, and the recorded calibration--test
group overlap is zero. Imaging conditions are not defect categories:
inpainting denotes an appearance artifact, whereas an unknown-defect label
denotes a physical-defect evaluation target. Defect novelty is evaluated
at the image level and is not added as a detector class.

Let $y_i=1$ denote an unknown defect and $y_i=0$ a normal image. We report
AUROC, AUPR, and FPR at 95\% true-positive rate, abbreviated FPR95. AUROC
and AUPR measure ranking over thresholds; FPR95 measures the normal-image
false-positive rate at the specified unknown-defect sensitivity. Higher
AUROC and AUPR and lower FPR95 are preferred. These metrics are fractions
in $[0,1]$. The positive prevalence in the 422-image tests is
$163/422\approx0.386$, which is the reference precision of random ranking.
All raw-score comparisons retain the score direction stored in the
evaluation files. A separately calibration-selected direction audit is
identified explicitly when discussed below.

The quality router has four observed target states: clean, rain, snow, and
inpainting. Its fifth configured output has no training target. The archived
router audit found clean-state predictions across the observed conditions.
Accordingly, the archived QCS and fusion results characterize the saved
score configurations and do not establish reliable automatic routing among
degradations. The three-branch results use separate coefficient vectors for
the three known conditions. They are reported as condition-specific fusion
analyses, rather than the performance of one globally fixed deployment rule.

The broader PLAMD-labelled evaluation is reported separately because its
37610-image population differs from the 422-image tests. We retain the
archived protocol names, \emph{Mixed PLAMD OOD} and \emph{Snow-held-out OOD}.
The latter is a protocol identifier; the saved image-level results alone
are insufficient to establish source-independent leave-one-defect-family-out
generalization. Neither protocol supplies an evaluation of unknown-defect
bounding-box localization.

\subsubsection{Computation and Statistical Reporting}
The detector benchmark uses a $640\times640$ input tensor, batch size one,
50 warm-up iterations, and 200 timed iterations on the selected CUDA device.
CUDA execution is synchronized before and after the timed loop. Latency and
FPS refer to the benchmarked model forward path; they exclude image loading
and preprocessing. Model-specific operations contained in the forward path
remain included. These numbers do not include an external DQCE adapter or
the complete score-calibration and rejection workflow.

Computation is estimated with THOP. We report profiled GFLOPs using the
explicit convention $\mathrm{FLOPs}=2\times\mathrm{MACs}$, where MACs denotes
multiply--accumulate operations. THOP does not necessarily
cover every custom or functional operator, so these are estimates for the
profiled forward path. Detector accuracy and the main unknown-defect
comparison are single-run archived results. DQCE ablations report the mean
and sample standard deviation across seeds 0, 1, and 2; differences are not
described as statistically significant without an accompanying test.

\subsection{Known-Component Detection and Model Cost}
\label{subsec:detection_results}

Table~\ref{tab:detection_complexity} combines detection accuracy with
parameter count, profiled computation, and measured runtime. QSD-DETR
achieves mAP$_{50}$ of 88.18\% and mAP$_{50:95}$ of 72.05\%, compared with
86.22\% and 70.59\% for RT-DETR. The corresponding observed gains are 1.97
and 1.47 percentage points. At the stricter IoU threshold of 0.75, the
observed gain is 1.13 percentage points. These results show that the evaluated
QSD-DETR checkpoint retains known-component detection while supporting the
additional image-level analysis.

\begin{table*}[!t]
  \centering
  \caption{Known-component detection accuracy and model-forward cost on the
  2626-image validation split. Detection metrics are percentages.
  GFLOPs$^\dagger$ are twice the archived THOP MAC counts and represent
  profiled operations. Runtime is for the benchmarked detector path, not the
  complete anomaly-assessment pipeline. ``--'' denotes an unavailable
  archived metric.}
  \label{tab:detection_complexity}
  {\QSDexpTableStyle
  \begin{tabular*}{\textwidth}{@{\extracolsep{\fill}}lrrrrrrrr@{}}
    \toprule
    Model & mAP$_{50}$ & mAP$_{75}$ & mAP$_{50:95}$ & mAP$_{50:75}$
    & Params (M) & GFLOPs$^\dagger$ & Latency (ms) & FPS \\
    \midrule
    YOLOv8m      & 88.40 & --    & 72.00 & --    & 25.867 & 79.120  & 5.774  & 173.18 \\
    YOLO11m      & 87.70 & --    & 70.60 & --    & 20.067 & 68.262  & 6.548  & 152.71 \\
    YOLOv12m     & 86.10 & --    & 69.10 & --    & 21.724 & 69.235  & 21.349 & 46.84 \\
    Faster R-CNN & 79.22 & 63.68 & 58.45 & --    & 41.439 & 182.324 & 16.355 & 61.14 \\
    RT-DETR      & 86.22 & 72.77 & 70.59 & 80.40 & 42.764 & 135.876 & 14.191 & 70.47 \\
    Ours         & 88.18 & 73.90 & 72.05 & 82.50 & 43.062 & 136.034 & 15.927 & 62.79 \\
    \bottomrule
  \end{tabular*}\par
  \smallskip
  \parbox{\textwidth}{YOLO accuracy is retained at the precision available in
  its archived validation logs. Faster R-CNN accuracy is from the supplied
  final validation log. RT-DETR and Ours use the recovered evaluation caches.
  Parameter counts and runtime use the saved benchmark checkpoints.}
  }
\end{table*}

YOLOv8m has the highest reported mAP$_{50}$, at 88.40\%, and the highest
measured detector throughput. Its mAP$_{50:95}$ of 72.00\% is close to the
72.05\% of QSD-DETR; the available precision and single-run protocol do not
support a claim of meaningful superiority between these two entries.
Faster R-CNN obtains lower accuracy in the supplied validation report and
has the largest profiled computation among the compared forward paths.

Figure~\ref{fig:detection_visualization} provides the qualitative
counterpart of Table~\ref{tab:detection_complexity}. It compares the same
four validation images across all six detectors, together with ground truth.
The cases cover a relatively easy scene, a crowded arrangement, small
components, and overlapping components. Identical image crops, class colors,
and display thresholds permit inspection of missed detections, false
positives, and localization differences. These selected examples illustrate
individual outcomes; they do not replace the aggregate validation results.

\begin{figure*}[!t]
  \centering
  \QSDexpArtwork{figure/figure6/fig6_detection_comparison.png}
    {figure/figure7/fig7_open_world_comparison.png}{2.5in}
  \caption{Qualitative known-component detection comparison. Columns show
  the same easy, crowded, small-object, and overlap examples; rows show
  ground truth, YOLOv8m, YOLO11m, YOLOv12m, Faster R-CNN, RT-DETR, and Ours.
  All predictions are displayed at confidence 0.25 with consistent class
  colors. Enlarged regions use the same coordinates for every model.
  The displayed boxes denote component categories, not unknown-defect
  localization.}
  \label{fig:detection_visualization}
\end{figure*}

Relative to RT-DETR, the benchmarked QSD-DETR checkpoint adds 0.299 million
parameters before rounding the table entries and approximately 0.158 profiled
GFLOPs, equivalent to 0.079 GMACs. Measured latency increases by 1.736 ms,
from 14.191 to 15.927 ms, while throughput decreases from 70.47 to 62.79 FPS.
This comparison characterizes the cost of the benchmarked checkpoint; the
complete defect-assessment workflow is considered separately in
Section~\ref{subsec:runtime_analysis}.

\subsection{Unknown-Defect Scoring Under Image Degradation}
\label{subsec:unknown_scoring}

Table~\ref{tab:unknown_methods} compares unknown-defect scoring methods
under the three imaging conditions. MSP is the RT-DETR
confidence-based score used in the baseline comparison. The additional
methods are energy scoring, Gaussian feature scoring with 90th-percentile
aggregation, Deep kNN, ViM, and PatchCore. The proposed configurations
separate DQCE-only scoring, the archived QCS+semantic score, and
condition-specific three-branch fusion. These suffixes identify different
score configurations of QSD-DETR with distinct inference rules.

\begin{table*}[!t]
  \centering
  \caption{Image-level unknown-defect comparison under rain, snow,
  and inpainting. Each condition contains 259 normal and 163 unknown-defect
  images. Entries are seed-0 results with the archived raw-score direction.
  AUROC and AUPR are higher-is-better; FPR95 is lower-is-better.
  The three-branch row uses condition-specific coefficients.}
  \label{tab:unknown_methods}
  {\QSDexpTableStyle
  \begin{tabular*}{\textwidth}{@{\extracolsep{\fill}}lrrrrrrrrr@{}}
    \toprule
    & \multicolumn{3}{c}{Rain} & \multicolumn{3}{c}{Snow}
    & \multicolumn{3}{c}{Inpainting} \\
    \cmidrule(lr){2-4}\cmidrule(lr){5-7}\cmidrule(lr){8-10}
    Method & AUROC & AUPR & FPR95 & AUROC & AUPR & FPR95
    & AUROC & AUPR & FPR95 \\
    \midrule
    RT-DETR (MSP) & 0.374 & 0.311 & 0.973 & 0.474 & 0.392 & 0.977 & 0.429 & 0.328 & 0.965 \\
    Energy        & 0.413 & 0.322 & 0.981 & 0.470 & 0.356 & 0.961 & 0.472 & 0.366 & 0.934 \\
    Gaussian (q90)& 0.669 & 0.544 & 0.915 & 0.558 & 0.437 & 0.938 & 0.544 & 0.434 & 0.954 \\
    Deep kNN      & 0.620 & 0.509 & 0.903 & 0.543 & 0.418 & 0.900 & 0.567 & 0.442 & 0.969 \\
    ViM           & 0.412 & 0.318 & 0.973 & 0.453 & 0.339 & 0.981 & 0.474 & 0.352 & 0.950 \\
    PatchCore     & \textbf{0.845} & \textbf{0.746} & \textbf{0.537}
                  & \textbf{0.629} & 0.471 & \textbf{0.822} & 0.356 & 0.301 & 0.969 \\
    \midrule
    Ours (DQCE)   & 0.703 & 0.569 & 0.710 & 0.591 & 0.468 & 0.853
                  & \textbf{0.748} & \textbf{0.647} & \textbf{0.695} \\
    Ours (QCS+semantic)
                  & 0.730 & 0.586 & 0.737 & 0.566 & 0.426 & 0.846 & 0.563 & 0.416 & 0.876 \\
    Ours (three-branch)
                  & 0.763 & 0.637 & 0.683 & 0.577 & \textbf{0.482} & 0.942 & 0.638 & 0.552 & 0.888 \\
    \bottomrule
  \end{tabular*}}
\end{table*}

The results show substantial condition dependence. Under rain, PatchCore
achieves the strongest reported AUROC, AUPR, and FPR95, at
0.845, 0.746, and 0.537. Among the proposed raw-score configurations,
three-branch fusion improves AUROC from 0.703 for DQCE to 0.763 and AUPR
from 0.569 to 0.637. Both exceed the archived RT-DETR confidence score,
although neither exceeds PatchCore in this condition.

Under inpainting, DQCE achieves AUROC/AUPR/FPR95 of 0.748/0.647/0.695,
the strongest values in the raw-score table. Its advantage over the
three-branch result, 0.638/0.552/0.888, shows that adding branches does not
automatically improve the final ranking. Snow is more difficult: PatchCore
has the highest AUROC and lowest FPR95, while three-branch fusion has the
highest AUPR of 0.482. The latter nevertheless has FPR95 of 0.942, so its
precision--recall advantage does not imply a useful high-sensitivity
operating point.

Figure~\ref{fig:condition_scores} complements the table with normal and
unknown score distributions for RT-DETR, PatchCore, and the proposed score
configurations. The class-conditional scores permit inspection of score
overlap, reversed ordering, and compressed score ranges beyond the aggregate
metrics. No defect localization
claim is inferred from these image-level scores.

\begin{figure*}[!t]
  \centering
  \QSDexpArtwork{figure/figure7/fig7_unknown_score_distributions.png}
    {figure/figure5/fig5_condition_comparison.png}{2.5in}
  \caption{Image-level score distributions under rain, snow, and inpainting
  on the condition-controlled test sets. Within each method--condition panel,
  normal and unknown-defect images are shown separately. RT-DETR, PatchCore,
  Ours (DQCE), Ours (QCS+semantic), and Ours (three-branch) use the same
  image identities as Table~\ref{tab:unknown_methods}. The raw-score direction
  is preserved. Method-specific score axes are stated explicitly because
  the scores are not calibrated probabilities on a shared scale.}
  \label{fig:condition_scores}
\end{figure*}

Score direction is a material part of this comparison. In a separate
calibration-based audit, the selected PatchCore direction is positive for
rain and snow and negative for inpainting. Across three seeds, its
direction-audited inpainting AUROC is $0.642\pm0.004$, AUPR is
$0.504\pm0.005$, and FPR95 is $0.816\pm0.012$. The corresponding DQCE
results are $0.717\pm0.027$, $0.643\pm0.029$, and $0.856\pm0.139$.
DQCE has higher mean AUROC and AUPR in this audited comparison, while
PatchCore has lower mean FPR95. This qualification prevents the raw
PatchCore inpainting result from being interpreted as an unconditional
failure of the method. Raw-direction and calibration-selected comparisons
are distinct protocols.

\subsection{Ablation of the DQCE Training Objective}
\label{subsec:dqce_ablation_results}

This experiment tests the contribution of paired consistency and
known-defect supervision while leaving the detector representation fixed.
The complete adapter uses the normal and known-defect binary losses, the
known-defect margin, and paired clean--degraded embedding consistency
described in the Method section. The \emph{w/o pair} configuration sets the
pair weight to zero. The \emph{w/o known supervision} configuration removes
known-defect samples and therefore removes both their binary loss and the
known-defect margin. It is not an ablation of the margin term alone.

\begin{table*}[!t]
  \centering
  \caption{DQCE objective ablation on the condition-controlled tests.
  Entries are the mean $\pm$ sample standard deviation over seeds 0, 1,
  and 2. ``w/o known supervision'' removes both the known-defect binary
  term and its margin. These repeated-run results are distinct from the
  seed-0 comparison in Table~\ref{tab:unknown_methods}.}
  \label{tab:dqce_ablation}
  {\QSDexpTableStyle
  \begin{tabular*}{\textwidth}{@{\extracolsep{\fill}}llccc@{}}
    \toprule
    Condition & Configuration & AUROC & AUPR & FPR95 \\
    \midrule
    Rain & Ours (full DQCE) & $0.667\pm0.031$ & $0.549\pm0.029$ & $0.822\pm0.097$ \\
         & Ours w/o pair & $0.681\pm0.021$ & $0.531\pm0.015$ & $0.810\pm0.080$ \\
         & Ours w/o known supervision & $0.713\pm0.006$ & $0.565\pm0.010$ & $0.678\pm0.012$ \\
    \midrule
    Snow & Ours (full DQCE) & $0.612\pm0.020$ & $0.513\pm0.042$ & $0.871\pm0.031$ \\
         & Ours w/o pair & $0.591\pm0.004$ & $0.462\pm0.004$ & $0.828\pm0.021$ \\
         & Ours w/o known supervision & $0.635\pm0.008$ & $0.476\pm0.008$ & $0.755\pm0.017$ \\
    \midrule
    Inpainting & Ours (full DQCE) & $0.717\pm0.027$ & $0.643\pm0.029$ & $0.856\pm0.139$ \\
         & Ours w/o pair & $0.705\pm0.019$ & $0.600\pm0.020$ & $0.699\pm0.038$ \\
         & Ours w/o known supervision & $0.684\pm0.011$ & $0.559\pm0.011$ & $0.797\pm0.045$ \\
    \bottomrule
  \end{tabular*}}
\end{table*}

Table~\ref{tab:dqce_ablation} shows that removing paired consistency lowers
mean AUPR in every condition: from 0.549 to 0.531 under rain, from 0.513 to
0.462 under snow, and from 0.643 to 0.600 under inpainting. AUROC and FPR95
do not improve uniformly with this term. In particular, removing it reduces
inpainting FPR95 from 0.856 to 0.699. The evidence therefore supports a
mean AUPR benefit in these runs, rather than improvement at every threshold
or on every metric.

Removing known-defect supervision also has mixed effects. The full adapter
has higher mean AUPR under snow and inpainting, but lower mean AUROC and
higher FPR95 under rain and snow. The saved logit statistics provide an
additional diagnostic: without known supervision, the within-run logit
standard deviation is approximately $6\times10^{-5}$, compared with values
of order one for the full adapter. Such score compression does not by itself
force AUROC to chance, since ranking is invariant to strictly monotonic
rescaling. It does, however, show why favorable ranking alone is insufficient
to establish well-separated or numerically stable energy evidence.

Figure~\ref{fig:dqce_ablation} therefore reports individual seed outcomes
alongside mean AUPR and a logit-dispersion diagnostic. This presentation adds
run-to-run variability and score-scale information to the exact aggregate
metrics in Table~\ref{tab:dqce_ablation}.

\begin{figure*}[!t]
  \centering
  \QSDexpArtwork{figure/figure8/fig8_dqce_ablation.png}
    {figure/figure6/fig6_dqce_ablation.png}{2.5in}
  \caption{Repeated-run DQCE ablation. Panels (a)--(c) show seed-wise AUPR
  and the mean $\pm$ sample standard deviation for rain, snow, and
  inpainting, respectively. Panel (d) shows the saved within-run logit
  standard deviation on a logarithmic axis. All configurations use seeds
  0, 1, and 2. The same seed bundle underlies
  Table~\ref{tab:dqce_ablation}; logit dispersion diagnoses score compression
  and is not a substitute for a ranking metric.}
  \label{fig:dqce_ablation}
\end{figure*}

\subsection{Broader Image-Level Open-World Evaluation}
\label{subsec:plamd_results}

The PLAMD-labelled protocols test the saved detector evidence on a larger
image population. Each archived evaluation contains 37610 images, including
24690 normal images and 12920 unknown-labelled images. These results are
not pooled with Table~\ref{tab:unknown_methods}: the sample population and
unknown prevalence differ. Table~\ref{tab:open_world} reports only
method--protocol pairs with a corresponding result artifact. In particular,
the semantic configuration is available for Snow-held-out OOD; no mixed
protocol value is inferred from that file.

\begin{table*}[!t]
  \centering
  \caption{Image-level results on the archived PLAMD-labelled protocols.
  Each protocol contains 24690 normal and 12920 unknown-labelled images.
  AUROC and AUPR are higher-is-better; FPR95 is lower-is-better.
  Comparisons are made within a protocol. These metrics do not evaluate
  unknown-defect bounding-box localization.}
  \label{tab:open_world}
  {\QSDexpTableStyle
  \begin{tabular*}{\textwidth}{@{\extracolsep{\fill}}llrrr@{}}
    \toprule
    Protocol & Method & AUROC & AUPR & FPR95 \\
    \midrule
    Mixed PLAMD OOD & RT-DETR & 0.398 & 0.312 & 1.000 \\
                    & Ours (QCS--CSR) & 0.609 & 0.498 & 0.760 \\
    \midrule
    Snow-held-out OOD & RT-DETR & 0.305 & 0.256 & 0.999 \\
                      & Ours (QCS--CSR) & 0.736 & 0.576 & 0.710 \\
                      & Ours (semantic configuration) & 0.616 & 0.390 & 0.888 \\
    \bottomrule
  \end{tabular*}}
\end{table*}

In Mixed PLAMD OOD, QCS--CSR increases AUROC from 0.398 to 0.609 and AUPR
from 0.312 to 0.498 relative to the archived RT-DETR score. FPR95 decreases
from 1.000 to 0.760. Under Snow-held-out OOD, the corresponding AUROC
increases from 0.305 to 0.736, while FPR95 decreases from 0.999 to 0.710.
The semantic configuration improves on RT-DETR in the latter protocol but
does not exceed QCS--CSR on the three reported metrics.

Figure~\ref{fig:open_world_scores} compares normal and unknown score
distributions for RT-DETR and QCS--CSR within each protocol. Its purpose is
to expose the score ordering behind Table~\ref{tab:open_world}.
The observed improvements support the
use of additional detector evidence in these archived image-level tests.
They do not isolate automatic quality routing, establish performance for
the full three-branch pipeline, or demonstrate localization of novel defects.

\begin{figure*}[!t]
  \centering
  \QSDexpArtwork{figure/figure9/fig9_open_world_scores.png}
    {figure/figure7/fig7_open_world_comparison.png}{2.5in}
  \caption{Normal and unknown image-level score distributions for RT-DETR
  and Ours (QCS--CSR), shown separately for Mixed PLAMD OOD and
  Snow-held-out OOD. Each panel preserves the saved score direction and
  identifies the method-specific score axis. Protocols are not pooled.
  The distributions complement the aggregate metrics in
  Table~\ref{tab:open_world}.}
  \label{fig:open_world_scores}
\end{figure*}

\subsection{Runtime and Practical Assessment Cost}
\label{subsec:runtime_analysis}

Table~\ref{tab:detection_complexity} reports detector-forward complexity
and runtime. Figure~\ref{fig:efficiency} relates detection accuracy to
latency and parameter count, and adds a separately timed image-level
assessment experiment. The latter is needed because detector FPS alone
does not measure the cost of loading images and running an external
anomaly-scoring adapter or memory bank.

In the saved assessment benchmark, the DQCE and PatchCore workflows process
the same 422-image evaluation set. The recorded mean wall times are 16.22 s
and 89.56 s, respectively, with run-to-run standard deviations of 0.85 s
and 3.22 s. These totals include model loading and image I/O. They correspond
to amortized rates of 26.02 and 4.71 images/s over this evaluation size,
with peak GPU memory of 2312 and 2466 MB in the recorded memory runs.
The measured wall-time ratio is approximately 5.52. These are workflow
measurements for the recorded implementations, rather than steady-state
detector FPS or a hardware-independent complexity claim.

\begin{figure*}[!t]
  \centering
  \QSDexpArtwork{figure/figure10/fig10_accuracy_efficiency.png}
    {figure/figure8/fig8_efficiency_comparison.png}{2.7in}
  \caption{Accuracy and practical assessment cost. (a) Detection
  mAP$_{50:95}$ versus detector-forward latency for the six models in
  Table~\ref{tab:detection_complexity}, with marker area indicating parameter
  count. (b) Repeated-run wall time for DQCE and PatchCore on 422 images,
  including model loading and image I/O; error bars denote the recorded
  run-to-run standard deviation. (c) Peak GPU memory in the corresponding
  memory runs. The detector-forward and full-workflow timing protocols
  are kept in separate panels.}
  \label{fig:efficiency}
\end{figure*}

\subsection{Empirical Calibration-Based Rejection}
\label{subsec:selective_results}

Ranking and thresholded rejection address different operating requirements.
For image $i$, let $r_i=1$ denote rejection as an unknown-defect candidate.
With $y_i=1$ for an unknown defect, acceptance coverage, selective risk,
and unknown recall are
\begin{align}
  \mathrm{Cov.}&=\frac{1}{N}\sum_i(1-r_i),\\
  \mathrm{Risk}&=\frac{\sum_i(1-r_i)y_i}{\sum_i(1-r_i)},\\
  \mathrm{Recall}_{\mathrm{unk}}&=\frac{\sum_i r_i y_i}{\sum_i y_i}.
  \label{eq:exp_selective_metrics}
\end{align}
Selective risk here is the unknown-image proportion among accepted images,
not a bounding-box error rate. Coverage is a retention measure and must be
interpreted jointly with risk and unknown recall. For a nonempty accepted
set, these quantities obey
\begin{equation}
 \mathrm{Risk}=\frac{(N_1/N)(1-\mathrm{Recall}_{\mathrm{unk}})}{\mathrm{Cov.}},
 \label{eq:exp_selective_relation}
\end{equation}
where $N_1$ is the number of unknown images. They therefore summarize
related aspects of the same acceptance decisions rather than independent
performance gains.

Table~\ref{tab:conformal} retains the archived operating points for three
score configurations: condition-calibrated QCS+semantic, the archived
unified QCS+semantic scoring variant, and condition-specific three-branch
fusion. The first configuration calibrates QCS and semantic evidence.
The term ``unified'' identifies the second
stored scoring variant and does not imply that all three conditions share
one coefficient vector or one threshold.

\begin{table*}[!t]
  \centering
  \caption{Empirical selective-rejection operating points on the 422-image
  condition-controlled tests. Cov. is acceptance coverage; Risk is the
  unknown-image fraction among accepted images; Recall is unknown recall.
  The first row uses condition-calibrated QCS+semantic scoring.
  Thresholds are fixed by their respective calibration artifacts.}
  \label{tab:conformal}
  {\QSDexpTableStyle
  \begin{tabular*}{\textwidth}{@{\extracolsep{\fill}}lrrrrrrrrr@{}}
    \toprule
    & \multicolumn{3}{c}{Rain} & \multicolumn{3}{c}{Snow}
    & \multicolumn{3}{c}{Inpainting} \\
    \cmidrule(lr){2-4}\cmidrule(lr){5-7}\cmidrule(lr){8-10}
    Configuration & Cov. & Risk & Recall & Cov. & Risk & Recall
    & Cov. & Risk & Recall \\
    \midrule
    \shortstack[l]{Ours (QCS+semantic,\\condition-calibrated)}
      & 0.893 & 0.347 & 0.196 & 0.929 & 0.367 & 0.117 & 0.884 & 0.359 & 0.178 \\
    \shortstack[l]{Ours (QCS+semantic,\\unified scoring)}
      & 0.877 & 0.338 & 0.233 & 0.917 & 0.372 & 0.117 & 0.882 & 0.358 & 0.184 \\
    Ours (three-branch)
      & 0.896 & 0.349 & 0.190 & 0.917 & 0.372 & 0.117 & 0.884 & 0.346 & 0.209 \\
    \bottomrule
  \end{tabular*}}
\end{table*}

At the three-branch operating point under rain, coverage is 0.896, selective
risk is 0.349, and unknown recall is 0.190. Under inpainting, the corresponding
values are 0.884, 0.346, and 0.209. Under snow, unknown recall is only 0.117
at coverage 0.917. Thus, most unknown images remain accepted at the selected
threshold, despite the gains in some ranking metrics. Figure~\ref{fig:conformal}
extends the single operating points to threshold-sweep risk--coverage and
unknown-recall--coverage curves. The sweep describes test-set behavior;
the table markers continue to use the stored calibration thresholds.

\begin{figure*}[!t]
  \centering
  \QSDexpArtwork{figure/figure11/fig11_selective_risk_coverage.png}
    {figure/figure9/fig9_conformal_rejection.png}{2.7in}
  \caption{Empirical selective-rejection tradeoffs under rain, snow, and
  inpainting. Risk--coverage and unknown-recall--coverage curves are obtained
  by sweeping the score threshold for evaluation. Markers identify the fixed
  calibration-selected operating points in Table~\ref{tab:conformal}.
  Configurations are Ours (QCS+semantic, condition-calibrated), Ours
  (QCS+semantic, unified scoring), and Ours (three-branch).
  The curves do not imply a distribution-free guarantee under condition
  shift, and no test-sweep threshold is substituted for the calibrated rule.}
  \label{fig:conformal}
\end{figure*}

These results are empirical selective-prediction measurements. The archived
workflow does not establish all independence and exchangeability conditions
needed for a finite-sample split-conformal guarantee, particularly when
score fitting and threshold calibration reuse observations or quality
conditions shift. Consequently, nominal calibration levels are not
interpreted as guaranteed test coverage or risk bounds.

\subsection{Condition-Specific Fusion Analysis}
\label{subsec:fusion_analysis}

The fusion analysis uses the standardized branch evidence defined in the
Method section. To make the experimental setting explicit, let
\begin{equation}
 S^{(d)}(\mathbf{x})=
 w_{\mathrm{qcs}}^{(d)}\widetilde{A}^{\mathrm{qcs}}(\mathbf{x})+
 w_{\mathrm{sem}}^{(d)}\widetilde{A}^{\mathrm{sem}}(\mathbf{x})+
 w_{\mathrm{dqce}}^{(d)}\widetilde{A}^{\mathrm{dqce}}(\mathbf{x}),
 \label{eq:exp_condition_fusion}
\end{equation}
where $d$ denotes the known evaluation condition. In QCS/semantic/DQCE
order, the archived coefficient vectors are $(0.25,0.45,0.30)$ for rain,
$(0.35,0.65,0.00)$ for snow, and $(0.05,0.10,0.85)$ for inpainting.
Each vector is nonnegative and sums to one. In particular, the snow
configuration assigns zero weight to DQCE, so it should not be described as
using three nonzero branch contributions.

Figure~\ref{fig:three_branch} visualizes these coefficients and image-level
rank correlations between the branch scores. The coefficients describe the
recorded combination, while the correlations permit inspection of whether
the branches order the same test images similarly. Neither statistic alone
establishes that a branch causes a gain. The actual fused ranking is
evaluated in Table~\ref{tab:unknown_methods} and its calibrated operating
point in Table~\ref{tab:conformal}.

\begin{figure*}[!t]
  \centering
  \QSDexpArtwork{figure/figure12/fig12_fusion_analysis.png}
    {figure/figure10/fig10_three_branch.png}{2.5in}
  \caption{Condition-specific evidence fusion analysis. (a) The archived
  QCS/semantic/DQCE coefficient vectors for rain, snow, and inpainting.
  (b)--(d) Pairwise Spearman rank correlations among the three image-level
  branch scores on the same test images within each condition. The
  correlation analysis uses the saved score orientations and is descriptive.
  Aggregate fused metrics and fixed-threshold rejection results are reported
  in Tables~\ref{tab:unknown_methods} and~\ref{tab:conformal}, respectively.}
  \label{fig:three_branch}
\end{figure*}

Under rain, the archived fusion improves ranking relative to DQCE alone,
whereas under inpainting the DQCE-only result is stronger. The snow result
also shows that an AUPR improvement can coexist with a poor FPR95. These
outcomes motivate condition-aware evidence analysis but do not establish
that the stored condition-specific coefficients constitute an automatically
routed, universally preferable model. A single fixed fusion artifact
evaluated on untouched conditions is a separate deployment question.

\subsection{Discussion}
\label{subsec:experimental_discussion}

The experiments support three bounded conclusions. First, the evaluated
QSD-DETR checkpoint retains competitive component-detection accuracy, with
an observed mAP$_{50:95}$ gain of 1.47 percentage points over the evaluated
RT-DETR checkpoint. Second, unknown-defect evidence is condition dependent:
PatchCore is strongest under rain in the raw-score comparison, DQCE is
strongest under inpainting, and snow remains difficult at high unknown
sensitivity. Third, paired consistency improves mean AUPR across the three
conditions in the DQCE ablation, while known-defect supervision has mixed
effects across ranking and operating-point metrics.

The remaining limitations concern what the evaluations measure. Selected
detection examples cannot establish dataset-wide localization improvements;
image-level novelty results cannot establish localization of unseen defects.
The archived clean-state routing behavior and condition-specific fusion
prevent a claim of validated automatic adaptation across conditions.
Detector-forward FPS and complete-assessment wall time measure different
execution scopes. Finally, empirical rejection under condition shift does
not provide a guaranteed conformal risk bound. Keeping these distinctions
explicit aligns the reported evidence with the scope of the proposed
framework.
